% GENERATED FILE. Edit canonical lessons, then run npm run generate:books.
% canonical-source-hash: fnv1a64:ab69be620b7e5f48
% canonical-lessons: TE-C271-sunnam, TE-C271-penku, TE-C271-plastikku, TE-C271-nurugu, TE-C271-cillu

\chapter{Lime, Roof tile, Plastic, Foam, Hole}
\label{ch:te-lime-roof-tile-plastic-foam-hole}
\hlchaptermodality{271}

\begin{chapteropening}
\textbf{By the end of this chapter:} \emph{I can say lime, a roof tile, plastic, foam and a hole.}

Complete the last lesson of chapter 271: I can say lime, a roof tile, plastic, foam and a hole.
\end{chapteropening}

\section[\texorpdfstring{\te{సున్నం} (sunnaṁ)}{సున్నం (sunnaṁ)}]{\te{సున్నం} (sunnaṁ) — lime, whitewash}

\label{lesson:TE-C271-sunnam}



\begin{quote}
Before the new one: say the Telugu for a mortar, then the Telugu for a pestle.
\end{quote}

\subsection*{You'll want to know: \te{సున్నం}}
\textbf{\te{సున్నం}} — \emph{sunnaṁ} — \textquotedblleft{}lime, whitewash\textquotedblright{}.

\begin{grammarlens}[title={What you've built}]
One more word for this chapter.
\end{grammarlens}

\subsection*{Your turn}
\begin{itemize}
\raggedright
  \item \emph{Say it:} \emph{sunnaṁ}
  \item \emph{Say it:} \emph{sunnaṁ}, once more
  \item \emph{Say it:} say \emph{sunnaṁ}
  \item \emph{Recall:} say the Telugu for a mortar, then the Telugu for a pestle, then say \emph{sunnaṁ} again
\end{itemize}

\begin{tcolorbox}[breakable,colback=teal!4,colframe=teal!35!black,title={Before you move on}]
What does \te{సున్నం} mean? (\textquotedblleft{}Lime, whitewash\textquotedblright{}.) Say it once more.
\end{tcolorbox}
\section[\texorpdfstring{\te{పెంకు} (peṅku)}{పెంకు (peṅku)}]{\te{పెంకు} (peṅku) — a roof tile}

\label{lesson:TE-C271-penku}



\begin{quote}
Before the new one: say the Telugu for a pestle, then the Telugu for lime.
\end{quote}

\subsection*{You'll want to know: \te{పెంకు}}
\textbf{\te{పెంకు}} — \emph{peṅku} — \textquotedblleft{}a roof tile\textquotedblright{}.

\begin{grammarlens}[title={What you've built}]
One more word for this chapter.
\end{grammarlens}

\subsection*{Your turn}
\begin{itemize}
\raggedright
  \item \emph{Say it:} \emph{peṅku}
  \item \emph{Say it:} \emph{peṅku}, once more
  \item \emph{Say it:} say \emph{peṅku}
  \item \emph{Recall:} say the Telugu for a pestle, then the Telugu for lime, then say \emph{peṅku} again
\end{itemize}

\begin{tcolorbox}[breakable,colback=teal!4,colframe=teal!35!black,title={Before you move on}]
What does \te{పెంకు} mean? (\textquotedblleft{}A roof tile\textquotedblright{}.) Say it once more.
\end{tcolorbox}
\section[\texorpdfstring{\te{ప్లాస్టిక్కు} (plāsṭikku)}{ప్లాస్టిక్కు (plāsṭikku)}]{\te{ప్లాస్టిక్కు} (plāsṭikku) — plastic}

\label{lesson:TE-C271-plastikku}



\begin{quote}
Before the new one: say the Telugu for lime, then the Telugu for a roof tile.
\end{quote}

\subsection*{You'll want to know: \te{ప్లాస్టిక్కు}}
\textbf{\te{ప్లాస్టిక్కు}} — \emph{plāsṭikku} — \textquotedblleft{}plastic\textquotedblright{}.

\begin{grammarlens}[title={What you've built}]
One more word for this chapter.
\end{grammarlens}

\subsection*{Your turn}
\begin{itemize}
\raggedright
  \item \emph{Say it:} \emph{plāsṭikku}
  \item \emph{Say it:} \emph{plāsṭikku}, once more
  \item \emph{Say it:} say \emph{plāsṭikku}
  \item \emph{Recall:} say the Telugu for lime, then the Telugu for a roof tile, then say \emph{plāsṭikku} again
\end{itemize}

\begin{tcolorbox}[breakable,colback=teal!4,colframe=teal!35!black,title={Before you move on}]
What does \te{ప్లాస్టిక్కు} mean? (\textquotedblleft{}Plastic\textquotedblright{}.) Say it once more.
\end{tcolorbox}
\section[\texorpdfstring{\te{నురుగు} (nurugu)}{నురుగు (nurugu)}]{\te{నురుగు} (nurugu) — foam, lather}

\label{lesson:TE-C271-nurugu}



\begin{quote}
Before the new one: say the Telugu for a roof tile, then the Telugu for plastic.
\end{quote}

\subsection*{You'll want to know: \te{నురుగు}}
\textbf{\te{నురుగు}} — \emph{nurugu} — \textquotedblleft{}foam, lather\textquotedblright{}.

\begin{grammarlens}[title={What you've built}]
One more word for this chapter.
\end{grammarlens}

\subsection*{Your turn}
\begin{itemize}
\raggedright
  \item \emph{Say it:} \emph{nurugu}
  \item \emph{Say it:} \emph{nurugu}, once more
  \item \emph{Say it:} say \emph{nurugu}
  \item \emph{Recall:} say the Telugu for a roof tile, then the Telugu for plastic, then say \emph{nurugu} again
\end{itemize}

\begin{tcolorbox}[breakable,colback=teal!4,colframe=teal!35!black,title={Before you move on}]
What does \te{నురుగు} mean? (\textquotedblleft{}Foam, lather\textquotedblright{}.) Say it once more.
\end{tcolorbox}
\section[\texorpdfstring{\te{చిల్లు} (cillu)}{చిల్లు (cillu)}]{\te{చిల్లు} (cillu) — a hole}

\label{lesson:TE-C271-cillu}



\begin{quote}
Before the new one: say the Telugu for plastic, then the Telugu for foam.
\end{quote}

\subsection*{You'll want to know: \te{చిల్లు}}
\textbf{\te{చిల్లు}} — \emph{cillu} — \textquotedblleft{}a hole\textquotedblright{}.

\begin{grammarlens}[title={What you've built}]
One more word for this chapter.
\end{grammarlens}

\subsection*{Your turn}
\begin{itemize}
\raggedright
  \item \emph{Say it:} \emph{cillu}
  \item \emph{Say it:} \emph{cillu}, once more
  \item \emph{Say it:} say \emph{cillu}
  \item \emph{Recall:} say the Telugu for plastic, then the Telugu for foam, then say \emph{cillu} again
\end{itemize}

\begin{tcolorbox}[breakable,colback=teal!4,colframe=teal!35!black,title={Before you move on}]
What does \te{చిల్లు} mean? (\textquotedblleft{}A hole\textquotedblright{}.) Say it once more.
\end{tcolorbox}
